\section{ How the community regards Franchi–Le Jan}
The reception is best described as respect for a mathematically sophisticated construction, accompanied by reservations about its physical interpretation.
The original work belongs primarily to stochastic differential geometry and probability theory, rather than to mainstream theoretical particle physics.
Its central idea is quite attractive:
\begin{itemize}
   \item The particle follows a timelike worldline.
   \item Its four-velocity remains on the relativistic mass shell.
   \item Randomness acts on the acceleration, with proper time as the evolution parameter.
   \item The construction is intrinsic to spacetime geometry and can be extended to curved Lorentzian manifolds.
    \item In schematic form, in flat spacetime,
\[dX^\mu_s=U^\mu_s\,ds,\qquad U_\mu U^\mu=c^2,\]
with stochastic evolution of \(U^\mu_s\) taking place in the tangent space of the mass hyperboloid.
\end{itemize}
This is a significant mathematical achievement. In particular, it avoids the elementary difficulty of ordinary Wiener motion in spacetime, whose trajectories need not respect the light cone.

The work has received sustained mathematical recognition. The publisher currently reports 61 Crossref citations for the 2007 paper. 
However, there is an important qualification.

\section{The main physical objection}
A particularly interesting assessment was given by F. Debbasch and C. Chevalier in their review Relativistic Stochastic Processes.
They distinguish several competing models:
\begin{table}[h]
    \begin{tabular}{p{0.4\textwidth}p{0.4\textwidth}}
        Model&	Basic physical interpretation\\
        Relativistic Ornstein–Uhlenbeck	&Particle interacting with a surrounding medium\\
        Franchi–Le Jan &	Intrinsic geometrical diffusion\\
        Dunkel–Hänggi &	Relativistic Brownian motion with a medium\\
        Oron–Horwitz &	Diffusion with an invariant evolution parameter\\
    \end{tabular}
\end{table}
Their criticism of Franchi–Le Jan is quite revealing.
The process contains no explicit surrounding medium, no friction mechanism and no conventional source of thermal fluctuations.
Consequently, it is difficult to interpret it as ordinary Brownian motion of a particle in a physical environment.
They suggest that one possible interpretation is an effective diffusion induced by quantum fluctuations of spacetime. But that interpretation is speculative, not an established physical explanation. 

I think this is an important distinction: the mathematical construction is not thereby invalidated. Rather, its physical ontology remains unsettled.

How the related theories compare:
\begin{description}
    \item [Franchi–Le Jan] The most geometrical approach. The mass hyperboloid and the geometry of spacetime determine the stochastic process. Particularly interesting for general relativity.
    \item [Dunkel–Hänggi] Attempts to construct a physically motivated relativistic Brownian motion, with interaction with a surrounding fluid. Their 2007 PRD paper proposes a continuous, non-Markovian relativistic diffusion in position space.
    \item [Debbasch and collaborators] Develop relativistic Ornstein–Uhlenbeck processes, kinetic equations, equilibrium distributions and fluctuation-dissipation relations. This is closer to conventional relativistic statistical physics.
\end{description}
One particularly useful mathematical development is Bailleul's 2010 paper, A stochastic approach to relativistic diffusions. It develops a unified treatment on Lorentzian manifolds and establishes an associated H-theorem. 

\section{ Why have these theories not become mainstream?}
I would identify five obstacles.
\begin{enumerate}
   \item Lack of a convincing microscopic mechanism.For ordinary Brownian motion we know what produces the noise: collisions with molecules. In relativistic stochastic mechanics, the origin of the noise is much less clear.
   \item The problem of the physical time parameter.Using proper time is mathematically natural, but an ensemble of particles has different proper times. Connecting individual stochastic trajectories to a common physical time and a measurable probability distribution requires additional work.
   \item Relativistic causality and Markovianity.There are substantial restrictions on relativistic Markov processes in spacetime. The Franchi–Le Jan construction avoids some of these difficulties by working on relativistic phase space, whereas the Dunkel–Hänggi approach explores non-Markovian processes.
   \item Connection with quantum mechanics.A relativistic diffusion is not automatically a quantum theory. One still has to explain how the stochastic dynamics produces the Klein–Gordon or Dirac equation, how spin arises, and how quantum interference is recovered.
   \item Experimental relevance.There is no widely accepted experimental prediction that distinguishes these approaches from conventional relativistic statistical mechanics or quantum theory.
\end{enumerate}
These are substantial obstacles, not merely sociological reasons for their limited popularity.

\section{ Relation to Nelson's mechanics}
This is where I think the subject becomes especially interesting.
I would distinguish three different levels:
\begin{table}[h]
    \begin{tabular}{p{0.4\textwidth}p{0.4\textwidth}}
        Level&	Status\\
        Relativistic stochastic differential geometry&	Well-established mathematical subject\\
        Relativistic stochastic particle dynamics&	Several serious, competing formulations\\
        Relativistic stochastic foundations of quantum mechanics&	Much less developed and not generally accepted\\      
    \end{tabular}
\end{table}
The distinction matters because Franchi–Le Jan solves a problem that is not identical to Nelson's problem.
Nelson starts with a stochastic process and seeks to recover quantum dynamics.
Franchi–Le Jan starts with relativistic geometry and constructs a stochastic process compatible with that geometry.
Combining these ideas is not straightforward. In particular, the Franchi–Le Jan diffusion does not itself produce a Dirac equation.

For your proposed synthesis of SHP and relativistic Nelson mechanics, I see three potentially fruitful directions:
\begin{enumerate}
   \item Geometric foundation: use the Franchi–Le Jan construction to define diffusion on the relativistic mass-shell bundle.
   \item Quantum dynamics: introduce forward and backward stochastic derivatives and investigate whether a suitable complex combination produces relativistic wave dynamics.
   \item Spin and gauge fields: replace ordinary stochastic transport by stochastic parallel transport in an appropriate spinor bundle, with the spin connection and electromagnetic gauge connection entering explicitly.
\end{enumerate}
The third direction is particularly interesting. It might provide a genuine geometric origin for spin rather than simply inserting spin into an already constructed stochastic equation.
There is an important caveat: one must preserve the mass-shell constraint and distinguish the stochastic process of a classical relativistic particle from the quantum dynamics of its wavefunction.

 References worth examining
I would start with these:
1. Franchi and Le Jan, Relativistic Diffusions and Schwarzschild Geometry (2007).
2. Debbasch and Chevalier, Relativistic Stochastic Processes (2006).
3. Dunkel, Talkner and Hänggi, Relativistic Diffusion Processes and Random Walk Models (2007).
4. Bailleul, A Stochastic Approach to Relativistic Diffusions (2010).
5. Franchi and Le Jan, Hyperbolic Dynamics and Brownian Motion: An Introduction (2012).

My overall conclusion: I would not dismiss Franchi–Le Jan as an unsuccessful physical theory. I would classify it as a serious geometrical construction whose physical interpretation remains open. Its mathematical framework could be useful even if its original interpretation as a model of physical Brownian motion turns out not to be correct.

For our purposes, I would actually regard it as a potentially valuable starting point for constructing relativistic stochastic mechanics on curved spacetime, provided we do not confuse that geometrical starting point with the derivation of quantum mechanics.

One question I think deserves particular attention is whether the Franchi–Le Jan diffusion can be modified to incorporate a genuine Nelson forward/backward structure while preserving the relativistic mass-shell condition. That is a much more interesting problem than simply writing down a relativistic Langevin equation.
